\documentclass[11pt]{article}

\usepackage{amssymb,bm,mathtools} % fonts
\usepackage{color,fullpage,graphicx,hyperref,parskip} % layout + media

\def\FT#1{\ensuremath{\mathcal{F}{#1}}}
\def\FTI#1{\ensuremath{\mathcal{F}^{-1}{#1}}}

\begin{document}

\title{Numerical solutions of wave equations}
\author{Manu Mannattil}
\maketitle

\section{1D wave equations}

We want to solve the 1D wave equation given by
%
\begin{equation}
  u_{tt} = u_{xx}\quad\text{with}\quad x \in [-1, 1].
\end{equation}
%

\subsection{FFT solution}

We Fourier transform $u$ in real space to compute $u_{xx}$:
%
\begin{equation}
  u_{xx}(x; t) = \FTI[q^{2}\hat{u}(t)],
\end{equation}
%
where $\hat{u}(t) = \FT[u(x; t)]$ is the Fourier transform of $u$.
Verlet integration in time then gives
%
\begin{equation}
  u(x; {t + \Delta t}) = 2u(x; t) - u(x; t - \Delta t) + (\Delta t)^{2}u_{xx}(x; t).
\end{equation}

\subsection{Chebyshev collocation}

If $D$ is the Chebyshev differentiation matrix, we compute $u_{xx}$ as
%
\begin{equation}
  u_{xx}(x; t) = D^{2}u(x; t).
\end{equation}
%
Verlet integration in time then gives
%
\begin{equation}
  u(x; {t + \Delta t}) = 2u(x; t) - u(x; t - \Delta t) + (\Delta t)^{2}u_{xx}(x; t).
\end{equation}
%
However, in general, $u(x; t + \Delta t)$ will not satisfy the imposed boundary conditions.
For that we need to replace the end-point values $u(\pm 1; t + \Delta t)$ so that the boundary conditions are satisfied.

\begin{itemize}
  \item In the case of Dirichlet boundary conditions, we just set $u(\pm 1; t + \Delta t ) = 0$ after each time step.
    This will obviously only work if the time step is very small.

  \item For Neumann BC, we need to solve for $u_{1} = u(-1; t)$ and $u_{N} = u(1, t)$ from the BC $Du = 0$ after each time step.
    To do so, let the entries of $D$ be
  %
  \begin{equation}
    D = 
    \begin{pmatrix}
      D_{1,1} & D_{1,2} & \cdots & D_{1,N-1} & D_{1,N}\\
      D_{2,1} & D_{2,2} & \cdots & D_{2,N-1} & D_{2,N}\\
      \vdots & \vdots & \vdots & \vdots & \vdots\\
      D_{N-1,1} & D_{N-1,2} & \cdots & D_{N-1,N-1} & D_{N-1,N}\\
      D_{N,1} & D_{N,2} & \cdots & D_{N,N-1} & d_{N,N}
    \end{pmatrix}
  \end{equation}
  %
  Then, the BC can be written as two linear equations in two unknowns $u_{1}$ and $u_{N}$:
  %
  \begin{equation}
    \begin{aligned}
      D_{1,1}\,{\color{red}u_{1}} + D_{1,2}\,u_{2} + \cdots + D_{1,N-1}\,u_{N-1} + D_{1,N}\,{\color{red}u_{1}} &= 0,\\
      D_{N,1}\,{\color{red}u_{1}} + D_{N,2}\,u_{2} + \cdots + D_{N,N-1}\,u_{N-1} + D_{N,N}\,{\color{red}u_{N}} &= 0.
    \end{aligned}
  \end{equation}
  %
  This can be rearranged and written in matrix form as
  %
  \begin{equation}
    \begin{pmatrix}
      D_{1,1} & D_{1,N}\\
      D_{N,1} & D_{N,N}\\
    \end{pmatrix}
    \begin{pmatrix}
      \color{red}
      u_{1}\\
      \color{red}
      u_{2}
    \end{pmatrix}
    =
    \begin{pmatrix}
      D_{1,2}\,u_{2} + \cdots + D_{1,N-1}\,u_{N-1}\\
      D_{N,2}\,u_{2} + \cdots + D_{N,N-1}\,u_{N-1}
    \end{pmatrix}.
  \end{equation}
  %
  So to impose Neumann BC, we solve the above linear system after each time step.

  \item Absorbing or open boundary conditions.
    The wave equation can be factorized as
    %
    \begin{equation}
      \left(\partial_{t}^{2} - \partial_{x}^{2}\right)u = \left(\partial_{t} - \partial_{x}\right)\left(\partial_{t} + \partial_{x}\right)u = 0.
    \end{equation}
    %
    The equations $(\partial_{t} \pm \partial_{x})u = 0$ represent one-way waves moving in $\pm x$ directions.
    If we impose
    %
    \begin{equation}
      \begin{aligned}
        \partial_{t}u - \partial_{x}u &= 0 \quad\text{at}\quad x = -1,\\
        \partial_{t}u + \partial_{x}u &= 0 \quad\text{at}\quad x = 1,\\
      \end{aligned}
    \end{equation}
    %
    it will effectively kill waves that reenter the domain $[-1, 1]$ after reflection.
    Thus, we only use the full wave equation in updating the interior points $(u_{2}, u_{3}, \ldots, u_{N-1})$.  
    We update the end point values $u_{1}$ and $u_{N}$ using the BCs instead:
    %
    \begin{equation}
      \begin{aligned}
        u(-1; t + \Delta t) &= u(-1; t) + \Delta t\,u_{x}(-1; t),\\
        u(1; t + \Delta t) &= u(1; t) - \Delta t\,u_{x}(1; t).\\
      \end{aligned}
    \end{equation}
    %

\end{itemize}

\end{document}
